% Part: first-order-logic
% Chapter: introduction
% Section: first-order-logic

\documentclass[../../../include/open-logic-section]{subfiles}

\begin{document}

\olfileid[fr]{fol}{int}{fol}

\olsection{Logique du premier ordre}

Vous connaissez probablement déjà la logique du premier ordre grâce à
votre première initiation à la logique formelle.\footnote{À vrai dire,
nous le supposons plus ou moins !{} Si ce n'est pas le cas, vous pouvez
consulter un manuel plus élémentaire, tel que \emph{forall x}
\citep{Magnus2021}.} Vous la connaissez peut-être sous le nom de
« logique quantifiée » ou de « logique des prédicats ». La logique du
premier ordre est d'abord un langage formel. Elle possède donc un
vocabulaire déterminé, et ses expressions sont des chaînes formées à
partir de ce vocabulaire. Toutes les chaînes ne sont cependant pas
admises. Il existe plusieurs sortes d'expressions admises : des termes,
des !!{formula}s et des !!{sentence}s. Ce sont surtout les
!!{sentence}s de la logique du premier ordre qui nous intéressent : ils
fournissent un analogue formel des phrases déclaratives de la langue
courante; à leur sujet, nous pouvons poser les questions qui intéressent
habituellement les logiciens. Par exemple :
\begin{itemize}
    \item L'!!{sentence}~$!B$ découle-t-il logiquement de~$!A$ ?
    \item L'!!{sentence}~$!A$ est-il logiquement vrai, logiquement faux
      ou contingent ?
    \item Les !!{sentence}s~$!A$ et~$!B$ sont-ils équivalents ?
\end{itemize}

Ces questions portent avant tout sur la « signification » des
!!{sentence}s de la logique du premier ordre. Un philosophe analyserait
par exemple la question de savoir si $!B$ découle logiquement de~$!A$ de
la manière suivante : existe-t-il un cas où $!A$ est vrai et~$!B$ faux
($!B$ ne découle pas de~$!A$), ou bien tout cas qui rend $!A$ vrai
rend-il aussi~$!B$ vrai ($!B$ découle alors de~$!A$) ? Mais on ne nous a
pas encore dit ce qu'est un « cas » : c'est précisément le rôle de la
\emph{sémantique}. La sémantique de la logique du premier ordre fournit
un modèle mathématiquement précis de l'idée intuitive de « cas » et,
point essentiel, de ce que signifie pour !!a{sentence}~$!A$ d'être
\emph{vrai dans} un cas. Nous appellerons !!a{structure} le modèle
mathématiquement précis que nous allons élaborer. La relation qui rend
précise l'expression « vrai dans » s'appelle la relation de
\emph{satisfaction}. Nous définirons donc « $!A$ est satisfait
dans~$\Struct{M}$ » (en symboles : $\Sat{M}{!A}$) pour des
!!{sentence}s~$!A$ et des !!{structure}s~$\Struct{M}$. Cela fait, nous
pourrons également définir avec précision d'autres termes sémantiques,
comme « découle de » ou « est logiquement vrai ». Ces définitions
permettront de décider, toujours avec une précision mathématique, si,
par exemple,
\[
  \lforall[x][(!A(x) \lif !B(x))],\,
  \lexists[x][!A(x)] \Entails \lexists[x][!B(x)].
\]
La réponse sera évidemment affirmative. Si vous avez déjà appris à
symboliser des phrases de la langue courante en logique du premier
ordre, vous reconnaîtrez ici, par exemple, la symbolisation de
« Toutes les fourmis sont des insectes ; il existe des fourmis ; donc
il existe des insectes ». Cet argument est manifestement valide : notre
modèle mathématique de « découle de » pour le langage formel doit donc
donner la même réponse.

Vous vous souvenez probablement aussi, depuis votre première initiation
à la logique formelle, qu'il existe des \emph{!!{derivation}s}. Si vous
avez suivi un premier cours de logique formelle, on vous aura fait
chercher de telles !!{derivation}s, peut-être même une
!!{derivation} montrant que la conséquence logique ci-dessus vaut. Il
existe de nombreuses manières de présenter des !!{derivation}s : vous
avez peut-être pratiqué ce qu'on appelle la « déduction naturelle » ou
les « arbres de vérité », mais il en existe bien d'autres. Les systèmes
de !!{derivation} visent à fournir des outils qui permettent de répondre
aux questions des logiciens énumérées plus haut. Ainsi, une
!!{derivation} de déduction naturelle dont
$\lforall[x][(!A(x) \lif !B(x))]$ et $\lexists[x][!A(x)]$ sont les
prémisses, et $\lexists[x][!B(x)]$ la conclusion (la dernière ligne),
\emph{établit} que $\lexists[x][!B(x)]$ découle logiquement de
$\lforall[x][(!A(x) \lif !B(x))]$ et de $\lexists[x][!A(x)]$.

Mais pourquoi ? À première vue, les systèmes de !!{derivation} n'ont
rien à voir avec la sémantique : produire une !!{derivation} formelle
consiste seulement à disposer des symboles d'une certaine manière,
conformément à des règles ; ni les « cas » ni l'expression « vrai dans »
n'y figurent. Le lien entre les systèmes de !!{derivation} et la
sémantique doit être établi par une étude métalogique. Il faut notamment
une démonstration mathématique du fait qu'une !!{derivation} formelle de
$\lexists[x][!B(x)]$ à partir des prémisses
$\lforall[x][(!A(x) \lif !B(x))]$ et $\lexists[x][!A(x)]$ est possible
si et seulement si $\lforall[x][(!A(x) \lif !B(x))]$ et
$\lexists[x][!A(x)]$ entraînent ensemble $\lexists[x][!B(x)]$. Avant
d'y parvenir, il faut toutefois accomplir un travail minutieux pour
définir correctement la syntaxe et la sémantique.

\begin{editorial}
Dans trois formules de la source, le second argument de
$\lforall[x]$ se ferme après la conséquence logique entière, au lieu de
se fermer après $(!A(x) \lif !B(x))$ ; une seconde fermeture apparaît
alors à la fin. Les limites d'arguments voulues sont rétablies dans les
trois occurrences.
\end{editorial}

\end{document}
